\documentclass[11pt,letterpaper]{article}
\usepackage[utf8]{inputenc}
\usepackage{geometry}
\usepackage{setspace}

\geometry{left=1.0in,right=1.0in,top=1.0in,bottom=1.0in}
\doublespacing

\setlength{\parindent}{0pt}
\setlength{\parskip}{12pt plus 2pt minus 2pt}

\begin{document}

\thispagestyle{empty}

\begin{center}
    \large\textbf{IN THE UNITED STATES PATENT AND TRADEMARK OFFICE} \\
\end{center}

\vspace{0.2cm}
\noindent\begin{tabular}{p{2.3in}p{4.0in}}
\textbf{In re Patent Application of:} & Reinaldo Javier Aguiar Marcano \\
\textbf{Application Number:} & 19/560,038 \\
\textbf{Confirmation Number:} & 2359 \\
\textbf{Filing Date:} & March 7, 2026 \\
\textbf{Group Art Unit:} & Pre-Exam / TC 2100 \\
\textbf{Examiner:} & ADAM DEMARAY \\
\textbf{Title of Invention:} & OPTICAL COMPUTING ARCHITECTURE AND SYSTEM \\
\end{tabular}

\vspace{0.6cm}
\begin{center}
    \textbf{\large SUBSTITUTE SPECIFICATION}
\end{center}

\vspace{0.4cm}

\begin{center}
    \textbf{\large OPTICAL COMPUTING ARCHITECTURE AND SYSTEM}
\end{center}

\section*{Background of the Invention}
For over half a century, the architecture of computers has been predicated on the movement of electrons over metallic traces and pins. While this paradigm has been incredibly successful, it is now facing fundamental physical limitations. The performance of modern processors has outpaced the ability of electrical buses to supply them with data. Key bottlenecks, such as the limited number of pins on a CPU package (the ``I/O wall'') and the latency of accessing DRAM over an electrical bus (the ``memory wall''), now constrain the real-world performance of entire systems. Issues of signal integrity, electromagnetic interference (EMI), power consumption, and thermal density further inhibit the scaling of high-speed electrical interconnects. A new architectural foundation is required to unlock the next generation of computing performance.

\section*{Summary of the Invention}
The present invention is a complete computer system architecture wherein the primary data pathways are optical rather than electrical. This ``Optical Computing Architecture'' leverages a unified physical layer based on massively parallel optical interconnects to create a holistic, high-bandwidth data fabric.

The system integrates several key optical components:
\begin{enumerate}
    \item \textbf{An Optical CPU Socket (``CPU Light Socket''):} The traditional pin grid array (PGA) or land grid array (LGA) of a CPU is replaced with a two-dimensional array of photon emitters and/or sensors, which optically couples the CPU to the motherboard.
    \item \textbf{An Optical Memory Bus:} The electrical DIMM slots and memory module connectors are replaced with optical interfaces, allowing for petabit-per-second data transfer rates between the CPU and main memory.
    \item \textbf{Optical Peripheral Interconnects:} External and internal peripheral connections (e.g., for GPUs, storage, or inter-node communication) are implemented using the same optical data cable technology, creating a unified I/O system.
\end{enumerate}

By replacing all major data buses with a common optical framework, the architecture achieves synergistic benefits beyond the improvement of any single component. The elimination of electrical bottlenecks results in a system with an aggregate bandwidth of multiple petabits per second. This enables processors to operate at their full potential, fundamentally changing the design considerations for software and algorithms in data-intensive fields like AI and HPC. The architecture also allows for greater physical density, reduced power consumption, and enhanced security.

\section*{Statement of License Grant to the United States Government}
The inventor, Reinaldo Javier Aguiar Marcano, hereby grants to the United States Government and its agencies a perpetual, royalty-free, non-exclusive license to practice and use the invention described herein for any and all governmental and non-commercial purposes.

\section*{Brief Description of the Drawings}
\noindent\textbf{FIG. 1} is a conceptual block diagram of the Optical Computing Architecture, showing the major components interconnected by optical data pathways.

\noindent\textbf{FIG. 2} is a diagram contrasting the bottlenecks in a traditional electrical computer architecture with the unimpeded data flow in the Optical Computing Architecture.

\noindent\textbf{FIG. 3} is a conceptual diagram illustrating how the Optical Computing Architecture enables high-density 3D chip stacking with optical inter-layer communication.

\section*{Detailed Description of the Invention}
Referring to \textbf{FIG. 1}, the Optical Computing Architecture is shown in a system-level block diagram. The central component is a motherboard or substrate (101). A processor (CPU/GPU) (102) is mounted to the motherboard not with electrical pins, but with an optical interface, the ``CPU Light Socket'' (103), comprising mated arrays of photon emitters and sensors. This socket provides the primary data link between the processor and the rest of the system.

The motherboard (101) also contains one or more optical memory slots, which connect to optical memory modules (104) via an Optical Memory Bus (105). Further, the motherboard includes optical peripheral ports (106), which connect to external devices (107) or other computer nodes via an external optical cable (108). All primary data pathways (103, 105, 106, 108) are massively parallel optical links, creating a unified, high-bandwidth data fabric for the entire system.

\textbf{FIG. 2} contrasts the proposed architecture with a conventional one. In a traditional electrical architecture (200), the CPU (201) is constrained by a ``CPU I/O Wall'' (202) due to a limited number of physical pins. It is further constrained by a ``Memory Wall'' (203) due to the limited bandwidth of the electrical memory bus. Data flow is therefore restricted at multiple points. In the Optical Computing Architecture (210), the CPU (211) is coupled to memory (213) and other components via ultra-high-bandwidth optical pathways (212, 214). The bottlenecks are removed, allowing for unimpeded, high-volume data flow throughout the entire system.

\textbf{FIG. 3} illustrates a key advantage of the architecture for future semiconductor design. A 3D stacked chip system (300) is shown, comprising multiple layers of silicon dies (301, 302, 303), which could be a mix of logic, memory, and specialized accelerator dies. In conventional 3D stacking, inter-die communication relies on thousands of microscopic electrical connections called Through-Silicon Vias (TSVs), which are complex and costly to manufacture. In the Optical Computing Architecture, this is replaced by optical interfaces. The top of one die (301) has an array of photon emitters (304), which communicates directly with a corresponding array of photon sensors (305) on the bottom of the die above it (302). This creates an ``Optical Inter-Layer Bus'' (306), allowing for ultra-high-bandwidth communication between the layers of the chip stack, enabling a new level of integration and performance.

\vspace{0.8cm}
\begin{center}
    \textbf{\large INVENTOR STATEMENT UNDER 37 CFR \S{} 1.125(b)}
\end{center}

Pursuant to 37 CFR \S{} 1.125(b), the undersigned Applicant hereby certifies that this substitute specification includes no new matter.

\vspace{0.6cm}
\begin{singlespace}
\noindent Respectfully submitted, \par\vspace{1.0cm}
\underline{\textbf{/Reinaldo Javier Aguiar Marcano/}} \\
\textbf{Reinaldo Javier Aguiar Marcano} \\
Pro Se Applicant \\
Mailing Address: 2302 Britton Ridge Drive, Katy, Texas 77494 \\
Telephone: (202) 599-3979 \\
Email: \texttt{rj.aguiar@gmail.com} (USPTO Account) \\
Email: \texttt{rj.aguiar.m@gmail.com} \\
Email: \texttt{reinaldo.aguiar.marcano@isbn9781971481296.com} \\
Date: October 10, 2026
\end{singlespace}

\end{document}
